%=========================================================
% Chapter: Results and Analysis
%=========================================================

\subsection{Pareto Front Comparison}

The core output of the multi-objective optimisation is the Pareto
front, illustrating the trade-off vectors that represent
non-dominated station-deployment configurations across the four
objectives. Both NSGA-II and MOPSO were run independently for
500~generations / iterations with identical evaluation budgets,
and their Pareto fronts are compared in
Figure~\ref{fig:pareto_comparison}.

\begin{figure}[H]
\centering
\includegraphics[width=0.95\linewidth]{image.png}
\caption{Pareto front comparison: NSGA-II vs.\ MOPSO across all
         four objective pairs (Cost vs.\ Coverage, Cost vs.\ Net
         ROI, Coverage vs.\ Net ROI, Net ROI vs.\ Mean Wait Time).}
\label{fig:pareto_comparison}
\end{figure}

\textbf{Interpretation of the Pareto front comparison:}
\begin{itemize}
    \item In the Cost vs.\ Coverage plane, both algorithms produce
          well-defined frontiers spanning the full coverage range.
          NSGA-II explores a wider coverage band (down to
          $\sim$46\%) by including very small low-cost layouts,
          while MOPSO concentrates its archive in the
          medium-to-high coverage band (76--100\%).
    \item In the Cost vs.\ Net ROI plane, both algorithms produce
          comparable maximum ROI ($\sim$4.2--4.3) at low cost,
          reflecting their similar ability to discover efficient
          small-network configurations. Such low-cost solutions,
          however, correspond to lower spatial coverage.
    \item In the Coverage vs.\ Net ROI plane, the classic
          efficiency trade-off is visible: high-coverage solutions
          sacrifice per-unit profitability, while low-coverage
          solutions cluster in the higher-ROI quadrant.
    \item The Cost vs.\ Coverage curve becomes notably steeper in
          the high-coverage region: the average marginal cost is
          approximately $\text{\rupee}6$--$7$\,L per percentage
          point of coverage, but pushing coverage to the upper end
          of the front incurs sharply diminishing returns as
          low-density peripheral wards require high-cost,
          low-utilisation deployments.
\end{itemize}

\textbf{Key Insight:}
The two algorithms are complementary on this problem instance:
NSGA-II provides a wider coverage span and a more uniformly spread
front, while MOPSO marginally wins on hypervolume and generational
distance. For a public-procurement decision, presenting the
\emph{union} of the two non-dominated archives gives the planner
the richest possible menu of trade-offs.

%---------------------------------------------------------
\subsection{Individual Pareto Front Results}

\begin{figure}[H]
\centering
\includegraphics[width=1\linewidth]{image1.png}
\caption{NSGA-II individual Pareto front results across all four
         objective pairs. 80~non-dominated solutions; cost range
         $\text{\rupee}71$\,L--$\text{\rupee}411$\,L; coverage
         range 46.4\%--100\%.}
\label{fig:nsga2_pareto}
\end{figure}

\begin{figure}[H]
\centering
\includegraphics[width=1\linewidth]{image3.png}
\caption{MOPSO individual Pareto front results across all four
         objective pairs. 100~archive solutions; cost range
         $\text{\rupee}39$\,L--$\text{\rupee}344$\,L; coverage
         range 75.8\%--100\%.}
\label{fig:mopso_pareto}
\end{figure}

\textbf{Interpretation of individual Pareto results:}
\begin{itemize}
    \item NSGA-II produces 80 non-dominated solutions with a
          coverage range of 46.4\%--100\% and a cost range of
          $\text{\rupee}71$\,L--$\text{\rupee}411$\,L. The
          per-solution Net ROI ranges from 1.85 to 4.15, reflecting
          genuine trade-off diversity across deployment scales.
    \item MOPSO produces 100 archive solutions with a coverage
          range of 75.8\%--100\% and a cost range of
          $\text{\rupee}39$\,L--$\text{\rupee}344$\,L. The
          per-solution Net ROI ranges from 2.12 to 4.33; the
          tighter low-end coverage is a consequence of MOPSO's
          archive-cap policy, which favours mid-to-high-budget
          regions.
    \item In the Net ROI vs.\ Mean Wait Time plane, all
          Pareto-optimal solutions from both algorithms cluster
          well below the 20-minute user-tolerance threshold; the
          maximum mean wait observed across either front is
          $\sim$5.6\,min, and no Pareto-optimal layout violates
          the $P_{\text{block}} > 10\%$ blocking constraint or the
          $\rho \ge 1$ stability bound.
\end{itemize}

%---------------------------------------------------------
\subsection{Algorithm Performance Indicators}

To formally quantify solution quality, three standard
multi-objective performance indicators are computed over the
four-dimensional objective space. The reference point for the
hypervolume calculation is set at 5\% above the empirical nadir of
the combined fronts.

\begin{table}[H]
\centering
\caption{Algorithm performance indicators}
\label{tab:performance}
\begin{tabular}{lrrr}
\toprule
Indicator & NSGA-II & MOPSO & Better \\
\midrule
Hypervolume (HV, normalised)        & 0.9967  & 1.0434  & Higher $\uparrow$ \\
Generational Distance (GD)          & $9.04 \times 10^{4}$ & $9.00 \times 10^{3}$ & Lower $\downarrow$ \\
Spread ($\Delta$)                   & 0.750   & 0.792   & Lower $\downarrow$ \\
Median $W_q$ (min)                  & 1.45    & 1.79    & Lower $\downarrow$ \\
Max Coverage (\%)                   & 100.0   & 100.0   & Higher $\uparrow$ \\
Max Net ROI                         & 4.15    & 4.33    & Higher $\uparrow$ \\
\bottomrule
\end{tabular}
\end{table}

The two algorithms split the indicators rather than one strictly
dominating the other. MOPSO achieves a marginally higher
hypervolume and a substantially lower generational distance,
indicating that its archive lies closer to the union-reference
front. NSGA-II produces a more evenly spread front and a slightly
lower median waiting time. Both algorithms achieve full coverage
(100\%) on at least one Pareto solution and reach comparable
maximum Net ROI values. The gap between the two HV values relative
to a common reference indicates that running both algorithms and
taking the union of their archives is preferable to relying on
either one individually.

%---------------------------------------------------------
\subsection{Optimal Solution Analysis}

For deployment recommendation, a candidate solution is extracted
from the NSGA-II Pareto front by selecting the layout that
maximises spatial coverage subject to the $P_{\text{block}} \le
10\%$ feasibility constraint across all selected stations. The
union archive contains layouts that achieve full
(100\%) ward coverage at a cost in the range
$\text{\rupee}3.5$\,Cr--$\text{\rupee}4.1$\,Cr, with mean queue
waits well within the 20-minute service-level target.

\begin{table}[H]
\centering
\caption{Representative high-coverage solutions from the NSGA-II
         Pareto front}
\label{tab:optimal}
\begin{tabular}{lrr}
\toprule
Metric & Value range & Unit \\
\midrule
Coverage ($f_2$)              & 100.0     & \%        \\
Total CAPEX ($f_1$)            & 3.51--4.11 & Crore INR \\
Net ROI ($f_3$)                & 2.18--2.77 & --        \\
Mean Queue $W_q$ ($f_4$)        & 0.27--0.39 & minutes   \\
Mean Pareto wait (full front)   & 0.13--5.55 & minutes   \\
\bottomrule
\end{tabular}
\end{table}

\textbf{Interpretation of the recommended solutions:}
\begin{itemize}
    \item Full-coverage layouts in the NSGA-II archive deploy on
          the order of 50--60 of the 118 candidate slots, with
          medium-sized parking hubs concentrated in dense Central
          and East zones where demand density justifies higher
          CAPEX relative to revenue generation.
    \item All Pareto-optimal solutions satisfy the per-zone
          grid-capacity constraint of 500\,kW, which the
          optimisation encodes as a soft penalty on $f_1$.
    \item A net ROI on the order of 2--3$\times$ across the
          high-coverage region of the front confirms that
          public--private partnership (PPP) deployment of EV
          charging stations in tier-2 Indian cities is
          commercially viable provided that site selection is
          rigorously optimised; the higher-ROI tail
          ($\sim 4\times$) corresponds to smaller, lower-coverage
          layouts suitable for a phased rollout.
\end{itemize}

%---------------------------------------------------------
\subsection{Queue Simulation Results}

\begin{figure}[H]
\centering
\includegraphics[width=1\linewidth]{image4.png}
\caption{Queue waiting time distribution (mean $W_q$ per Pareto
         solution) for NSGA-II (left) and MOPSO (right). Red dashed
         line: 20-minute user-tolerance threshold. Yellow dash-dot
         line: median (NSGA-II: 1.45\,min; MOPSO: 1.79\,min).}
\label{fig:queue_dist}
\end{figure}

\textbf{Interpretation of the queue waiting time distribution:}
\begin{itemize}
    \item NSGA-II solutions have a median real wait time of
          1.45\,min, ranging from 0.26 to 4.37\,min, all
          comfortably below the 20-minute user-tolerance threshold.
    \item MOPSO solutions have a slightly higher median real wait
          of 1.79\,min, with the full range spanning
          0.13--5.55\,min. The slightly higher median is
          consistent with MOPSO's preference for cheaper, more
          highly utilised configurations.
    \item No Pareto-optimal solution from either algorithm
          violates the $P_{\text{block}} > 10\%$ blocking
          threshold or the $\rho \ge 1$ stability bound; layouts
          that would have triggered the 50- or 20-minute penalty
          terms in $f_4$ are dominated and never appear on the
          final fronts.
    \item The inclusion of $W = 8$ physical waiting bays
          ($K = c + 8$) was important: an Erlang-B (loss-only)
          model would have predicted a higher customer-rejection
          rate during peak periods, whereas the finite buffer
          successfully queues the majority of the demand surge.
\end{itemize}

\textbf{Key Insight:}
Both algorithms produce service-feasible Pareto fronts under
realistic peak-load conditions. The ToU-weighted $M/M/c/K$
treatment is the mechanism that keeps the search away from
queue-infeasible layouts, regardless of which metaheuristic drives
the outer loop.

%---------------------------------------------------------
\subsection{Convergence Analysis}

\begin{figure}[H]
\centering
\includegraphics[width=1\linewidth]{image5.png}
\caption{NSGA-II convergence --- Hypervolume (HV, normalised) per
         generation. Final HV $= 0.9967$. The fastest growth occurs
         in the first $\sim$100 generations, after which the curve
         enters a long plateau.}
\label{fig:nsga2_conv}
\end{figure}

\begin{figure}[H]
\centering
\includegraphics[width=1\linewidth]{image5.png}
\caption{MOPSO convergence --- Hypervolume (HV, normalised) per
         iteration. Final HV $= 1.0434$. The leader-following
         dynamics produce a steeper rise in the first 50
         iterations followed by a plateau.}
\label{fig:mopso_conv}
\end{figure}

\begin{figure}[H]
\centering
\includegraphics[width=1\linewidth]{image8.png}
\caption{MOPSO HV improvement-rate analysis (50-iteration
         windows). The verdict is CONVERGED, validating the early
         stopping heuristic.}
\label{fig:mopso_rate}
\end{figure}

\textbf{Interpretation of convergence behaviour:}
\begin{itemize}
    \item NSGA-II grows monotonically to a final normalised HV of
          $0.9967$. The smooth cumulative curve confirms that
          elitist selection effectively preserves good solutions
          across generations.
    \item MOPSO reaches a slightly higher final normalised HV of
          $1.0434$, consistent with the known characteristic of
          particle swarms to converge faster in early iterations
          due to leader-following dynamics.
    \item The MOPSO improvement-rate plot confirms near-zero HV
          gain in the late iterations, with a formal CONVERGED
          verdict, validating the early-stopping heuristic
          ($\Delta\text{HV} < 5 \times 10^{-4}$ over a
          40-iteration window).
    \item Despite MOPSO's marginally higher final HV, NSGA-II
          produces a more uniformly spread front (Spread
          $0.750$ vs.\ $0.792$), so the two algorithms remain
          complementary.
\end{itemize}

%---------------------------------------------------------
\subsection{Sensitivity Analysis}

\begin{figure}[H]
\centering
\includegraphics[width=1\linewidth]{image6.png}
\caption{Annual profit distribution across Pareto solutions for
         NSGA-II (cyan) and MOPSO (pink). Profits span
         $\text{\rupee}1.2$\,Cr--$\text{\rupee}11.6$\,Cr/yr across
         the union archive.}
\label{fig:profit_dist}
\end{figure}

\textbf{Sensitivity findings} (one-factor sweep over EV growth
rate, electricity tariff, install-cost factor, and service rate):
\begin{itemize}
    \item \textbf{EV growth rate:} Scaling the logistic adoption
          curve upward increases station utilisation and net ROI
          but pushes a small subset of stations toward the
          $P_{\text{block}} = 10\%$ feasibility threshold,
          indicating that a network sized for 2025 demand will
          require hardware expansion well before 2030.
    \item \textbf{Electricity tariff:} An increase in the
          commercial grid tariff depresses net ROI roughly
          linearly, assuming the operator absorbs the cost
          without passing it through to users. ToU pass-through
          mechanisms would mitigate this sensitivity.
    \item \textbf{Charger service rate:} Reducing the mean session
          duration from 30 to 20 minutes (e.g.\ moving from a
          slower AC charger to a 50\,kW DC fast charger on the
          same battery) drops the time-weighted $\overline{W}_q$
          substantially and shifts the entire Pareto front in the
          cost-of-coverage direction. Conversely, doubling the
          session duration violates the 20-minute service level
          for a non-trivial subset of stations.
    \item \textbf{Profit distribution:} The annual-profit
          distribution across the union Pareto archive spans
          $\text{\rupee}1.2$\,Cr--$\text{\rupee}11.6$\,Cr per
          year, with NSGA-II reaching the upper end of the tail.
\end{itemize}

%---------------------------------------------------------
\subsection{Scalability Benchmarks}

\begin{figure}[H]
\centering
\includegraphics[width=1\linewidth]{image7.png}
\caption{Runtime scalability for NSGA-II and MOPSO as the
         candidate-pool size $N$ increases. Both algorithms scale
         polynomially.}
\label{fig:scalability}
\end{figure}

\textbf{Interpretation of scalability results:}
\begin{itemize}
    \item NSGA-II's per-step runtime grows polynomially with the
          candidate-pool size, consistent with its
          $\mathcal{O}(M N^{2})$ non-dominated sorting complexity
          (where $M$ is the number of objectives).
    \item MOPSO's per-step runtime is dominated by archive
          management and crowding-distance recomputation and
          remains approximately flat across the tested range.
    \item Full convergence runtimes for the Indore instance
          ($n = 118$ after dynamic candidate generation) complete
          in under four minutes of wall-clock time on a commodity
          laptop, well within practical planning timescales.
    \item The polynomial scaling confirms that the Python-based
          framework is applicable to larger metropolitan instances
          (e.g.\ Mumbai, Delhi with $n = 200$+ candidates) without
          requiring high-performance computing resources.
\end{itemize}

%---------------------------------------------------------
\subsection{Overall Interpretation}

Across all experimental analyses, several consistent conclusions
emerge:

\begin{itemize}
    \item NSGA-II and MOPSO are complementary on this binary EVCS
          placement problem. MOPSO marginally wins on hypervolume
          and generational distance; NSGA-II wins on spread and
          covers a wider coverage range. Presenting the union of
          the two archives to the planner therefore dominates
          relying on either algorithm individually.
    \item The $M/M/c/K$ queueing model with ToU segmentation and
          the embedded blocking-probability penalty proved
          essential: it kept the multi-objective search away from
          queue-infeasible layouts, so that no Pareto-optimal
          solution from either algorithm violates the
          $P_{\text{block}} > 10\%$ or $\rho \ge 1$ bounds.
    \item The Net ROI objective (annual profit divided by capex)
          successfully decoupled financial and spatial objectives
          and produced genuine Pareto diversity. ROI values across
          the union archive range from 1.85 to 4.33, with the
          higher values corresponding to small high-utilisation
          layouts and the lower values to broad full-coverage
          layouts.
    \item Monte~Carlo demand averaging ($S = 100$ independent
          per-station scenarios) ensures that the recommended
          deployments are structurally robust to the full range
          of plausible demand trajectories rather than only the
          mean forecast.
\end{itemize}

Together, these results verify that the optimised EV charging
network is economically feasible, operationally stable under peak
demand, and deployment-ready for Indore's anticipated EV adoption
trajectory through 2030 and beyond.
